% Transcription — fonds Alexandre Grothendieck, shelfmark 80, batch 4 (pages 61-72).
% Established from the facsimile archives/batches/80/batch-04.pdf, cut from a
% copy of 80.pdf fetched through the project's own relay
% (https://grothendieck-relay.onrender.com/source/80.pdf), not uploaded by the
% user; PDF page k = archivists' page 60+k, no cover sheet. Per the skill
% transcribe-grothendieck. The folder is seventy-two pages in four batches
% (1-20, 21-40, 41-60, 61-72), transcribed in parallel; this is the fourth and
% last, twelve pages. Batch 1 and batch 2 (transcripts/80/batch-01.fr.tex,
% batch-02.fr.tex) were read for notation and not edited; batch 3
% (batch-03.fr.tex) appeared while this pass was being checked and was read
% and aligned with, not edited: its p. 60, headed « Système de Pseudodroites »
% (a map on the real projective plane, vertices of even order >= 4, X̃ the
% orientation cover), breaks off at « Notons que sur l'ens. des sommets se
% trouvant », which p. 61 here continues. Its ν_s (order of a vertex) matches
% the ν(s) read on p. 67 here.
%
% Pass: Opus 5.5 (claude-opus-5-5), 2026-09-26 — first pass, unchecked against the pages by a human.
%
% Dating and title. \dating is the folder's own date, copied verbatim from
% src/content/catalogue.ts; the group « Géométrie et topologie combinatoire »
% (69-89) is not added, since the folder has a date of its own. Title from the
% catalogue, without its final full stop, as in batches 1 and 2.
%
% Pages skipped: 71, a blank cream separator sheet (only show-through from the
% neighbouring leaf). Pages summarised: none. Pages transcribed from his hand:
% 61-66 (pencil, A4 portrait, fast drafting; 66 has six lines only), 67-70 and
% 72 (pencil, smaller sheets in landscape). P. 67 is cancelled by diagonal
% hooks over its item c) and is transcribed under \struck. The drawings of
% pp. 63 and 69 are described in one French \note each, with their labels.
%
% His pagination: none in this batch. His numbered runs: items a) b) c) on
% p. 64 (properties of the bijections φ_{s;D,D'}), items a) b) c) on p. 67
% (c) struck), taken up again as c), (d) (circled), e) on p. 68, (e) and (f)
% (circled) on p. 69, and referred to as « a) b) c) d) e) f) » by the « Th »
% of p. 70; « 1°) » on p. 69 has no « 2°) ». P. 72 is a list by number of
% lines (1, 2, 3, 4, 5, n ≥ 6 droites) with sub-items a) b) c).
%
% What the batch is about.
%   61-66  Continuing « Système de Pseudodroites » (p. 60, batch 3): for a
%          pseudo-line D of the map, the vertices on D form a combinatorial
%          polygon and π_0(X̃ \ D̃) ↔ orientations of that polygon; for a
%          vertex s not on D, Ω_s(X) (orientations of X at s) ↔ Ω(D). He asks
%          how to make the composite bijection Ω(D) ≃ Ω(D') explicit in terms
%          of the 1-complex K^1 and its pseudo-lines (pp. 61-64): through two
%          lines Δ, Δ' at s meeting D, D' in u, u' (v, v'), t = D ∩ D', with a
%          special case (s of order 4, two pencils) handled by auxiliary lines
%          D'', Γ, Γ' (figure p. 63). P. 64 axiomatises: K^1 a 1-complex
%          subdivided into pseudo-lines, bijections φ_{s;D,D'} : Ω(D) ≃ Ω(D')
%          with properties a) simple transitivity, b) independence of s along
%          an arc of Δ cut by x, x', c) (cut off). P. 65: uniqueness fails with
%          a monogon or bigon (a double pencil); otherwise Ω(s) := quotient of
%          the Ω(D) by the φ_{s;D,D'}, a 2-element set, and a local system on
%          K^1 via Ω(s) ≃ Ω(D) ≃ Ω(s'), independent of D by b). P. 66: to
%          rebuild the map one would still need a polygonal structure on the
%          arcs at each vertex.
%   67-70  Restart with oriented pseudo-lines: D⃗(s) the oriented lines through
%          s (= arcs from s), ρ_{D⃗,s} a circular permutation of D⃗(s) for s
%          not on D; axioms a)-f) (ρ_{-D⃗,s} = ρ_{D⃗,s}^{-1}; comparison along Δ;
%          (d) ρ_{D⃗,s}(Δ⃗_1) meets D⃗ at the vertex following u_1 among the
%          intersections of D⃗ with the lines through s; (e) with figure;
%          (f) ρ^n(Δ⃗) = -Δ⃗ for s of order 2n). P. 70 « Th »: without monogons
%          or bigons such a system is unique, and defines the local system of
%          the Ω(s) = {ρ_s, ρ_s^{-1}}; marginal: d) e) f) suffice.
%   72     A table of arrangements of 1 to 5 and n ≥ 6 lines, general or
%          concurrent position: whether « standard proj. » / « standard
%          comb. », and projective endomorphisms vs combinatorial automorphisms
%          (Id, 𝔖_4, extensions of 𝔖_3 by Z/2, groups of the hexagon, octagon,
%          pentagon).
%
% Notation fixed once for the batch: \vec{D} for his arrowed D (oriented
% pseudo-line); \vec{\mathcal{D}} for his arrowed cursive D (the set of
% oriented pseudo-lines) and \vec{\mathcal{D}}(s); \rho for the long curled ρ
% of pp. 67-70 (reading of the letter itself doubtful, said once in a note);
% \Omega(D), \Omega_s(X), \Omega(s) as he writes them; K^1 as K^{1};
% \mathfrak{S} for his gothic S, as in batches 1-2; « ps. droites » for
% pseudo-droites, as in batches 1-2; \mathbf{Z}/2 for his Z/2.
%
% Readings to check first: the first line of p. 61 and its diagonal margin;
% the lines of p. 61 bottom on Δ, Δ', u, u', v, v', t; the start of p. 62
% (\overrightarrow{tu}); p. 63 middle (« pourvu que ... par malheur ») where a
% phrase seems missing; the first line of p. 65 (« serait unique par ... »);
% the superscript of Δ⃗_2 on pp. 68-69; the margin of p. 70; the right column
% of p. 72 (3 droites a) and b)).

\documentclass[11pt,a4paper]{article}
% Le préambule commun à toutes les transcriptions du fonds Grothendieck.
%
% Il tient en une idée : **l'incertitude se compose, elle ne se résout pas.**
% Une transcription qui lisse un mot douteux a détruit la seule chose pour
% laquelle on la faisait — un lecteur doit pouvoir savoir, à chaque endroit,
% ce qui était sur le feuillet et ce qui a été ajouté. Les six macros qui
% suivent sont donc l'appareil critique, et non de la décoration.
%
% Les mêmes commandes sont reconnues par `scripts/render.mjs`, qui produit la
% vue de lecture du site. Une macro ajoutée ici doit l'être aussi là-bas,
% faute de quoi le rendu échouera bruyamment — ce qui est voulu : un rendu
% silencieusement faux est pire qu'un rendu absent.

\ProvidesPackage{grothendieck}[2026/08/08 Transcriptions du fonds Grothendieck]

\usepackage{amsmath,amssymb,amsthm}
\usepackage{mathrsfs}
\usepackage{stmaryrd}
% Les diagrammes commutatifs sont partout dans ce fonds : c'est la langue
% même de la géométrie algébrique de Grothendieck, et une transcription qui
% les rendrait en prose ne transcrirait rien.
\usepackage{tikz-cd}
% Français par défaut — c'est la langue des feuillets — mais l'anglais est
% chargé aussi : la traduction et le résumé sont des documents anglais, et la
% typographie française (espaces avant les deux-points) y serait une faute.
% Ces documents basculent par \selectlanguage{english} en tête de corps.
\usepackage[english,french]{babel}
\usepackage{fontspec}
\usepackage{geometry}
\geometry{a4paper,margin=25mm,marginparwidth=28mm}
\usepackage{marginnote}
\usepackage{xcolor}
% ulem, not soul. `soul`'s \st and \ul are built on a character-by-character
% scan that XeLaTeX's font machinery breaks: the document fails with an
% undefined control sequence rather than degrading. ulem does the same two jobs
% and is Unicode-engine safe. `normalem` stops it hijacking \emph, which the
% transcriptions use for Grothendieck's own emphasis.
\usepackage[normalem]{ulem}
\usepackage{hyperref}
\hypersetup{colorlinks=true,linkcolor=black,urlcolor=black,pdfborder={0 0 0}}

% --- Métadonnées du lot -----------------------------------------------------
% Reprises telles quelles par le script de rendu, qui les affiche en tête de
% la vue de lecture. Elles fixent ce que le fichier prétend couvrir : sans
% elles, une transcription flotte sans point d'attache dans un fonds de seize
% mille feuillets.

\newcommand{\folder}[1]{\gdef\grothFolder{#1}}
\newcommand{\batch}[1]{\gdef\grothBatch{#1}}
\newcommand{\leaves}[2]{\gdef\grothFirst{#1}\gdef\grothLast{#2}}
\newcommand{\foldertitle}[1]{\gdef\grothFoldertitle{#1}}
\gdef\grothFolder{?}\gdef\grothBatch{?}\gdef\grothFirst{?}\gdef\grothLast{?}\gdef\grothFoldertitle{}

% --- Le résumé --------------------------------------------------------------
%
% Il ouvre la lecture modernisée, sous le titre « Résumé », et il est composé
% en petit. Ce n'est pas un ornement : c'est le seul passage du document qui
% ne soit pas des mathématiques, et un lecteur doit pouvoir le reconnaître
% avant de l'avoir lu — puis le sauter s'il connaît déjà le sujet, ou s'y
% arrêter s'il ne le connaît pas. Le filet qui le suit marque la reprise du
% corps.

\newenvironment{resume}
  {\small\setlength{\parskip}{0.45em}}
  {\par\medskip\hrule\medskip}

% --- Mots-clés ---------------------------------------------------------------
%
% En anglais, en clôture du résumé de la lecture modernisée : le vocabulaire
% moderne sous lequel chercher ce que les feuillets construisent. Le manifeste
% du site (`npm run manifest`) les extrait pour étiqueter la cote entière —
% cette ligne est la source unique des tags, il n'y a pas de fichier à côté.
\newcommand{\keywords}[1]{\par\smallskip\noindent
  {\footnotesize\textsc{Keywords} --- \textit{#1}}\par}

% --- L'appareil critique ----------------------------------------------------

% Le numéro de feuillet, en marge. C'est l'ancre qui permet de régler un
% désaccord sur une formule en regardant *un* feuillet plutôt qu'une chemise
% de six cents. Le site s'en sert aussi pour tourner les pages du fac-similé
% quand on fait défiler la transcription.
\newcommand{\leaf}[1]{%
  \marginnote{\footnotesize\color{gray!70}\textsf{#1}}%
  \label{leaf:#1}\ignorespaces}

% Illisible : on ne devine pas. Un mot inventé qui se lit comme les autres est
% le pire résultat possible.
\newcommand{\ill}{\textcolor{red!60!black}{[\,\ldots\,]}}

% Lecture douteuse : le mot est donné, et signalé comme incertain.
\newcommand{\uncertain}[1]{\uline{#1}}

% Ajout de l'éditeur — une lettre manquante, un mot sous-entendu.
\newcommand{\add}[1]{\textcolor{blue!50!black}{[#1]}}

% Biffé par Grothendieck lui-même. Ce qu'il a rayé fait partie du document :
% une rature dit souvent où la pensée a changé de direction.
\newcommand{\struck}[1]{\sout{#1}}

% Note du transcripteur, sur l'état matériel du feuillet.
\newcommand{\note}[1]{\par\smallskip{\small\color{gray!60!black}\textit{[#1]}}\par\smallskip}

% Note marginale de Grothendieck, distincte du corps du texte.
\newcommand{\marginal}[1]{\par\smallskip\noindent\hspace{1em}%
  {\small\color{gray!60!black}#1}\par\smallskip}

% --- Titre ------------------------------------------------------------------

\AtBeginDocument{%
  \begin{center}
    {\large\bfseries Fonds Alexandre Grothendieck --- cote n\textsuperscript{o}~\grothFolder}\\[2pt]
    {\itshape\grothFoldertitle}\\[4pt]
    {\small feuillets \grothFirst--\grothLast\ (lot \grothBatch)}\\[2pt]
    {\footnotesize\color{gray!60!black}Transcription établie d'après le fac-similé numérisé
     par l'Université de Montpellier.\\
     Les lectures incertaines sont soulignées, les illisibles notées [\,\ldots\,],
     les ajouts entre crochets.}
  \end{center}
  \vspace{1em}}

\endinput


\folder{80}
\batch{4}
\pages{61}{72}
\dating{[à partir de 1981]}
\watermark{Édition de démonstration}
\foldertitle{Systèmes de pseudo-droites : notes manuscrites (s.d.)}

\begin{document}

\page{61}
\note{la page continue la phrase interrompue au bas de la page 60 : « Notons que sur l'ens.\ des sommets se trouvant »}
sur une même $D$ il y a une structure polygonale induite (cela a un sens \add{clair}, \uncertain{vu} \uncertain{que} \add{l'ens.} $S_D$ de ces sommets est $\geqslant 3$). (Les orientations de $D$ \uncertain{comme} \uncertain{on} \uncertain{sait}, l'ens.\ $\pi_0(\widetilde{X} \setminus \widetilde{D})$) \uncertain{sont} en corr.\ 1-1 avec l'ens des orientations de ce polygône combinatoire. \struck{(Si $s_D = 2$ ou $s_D = 1$}
\note{« clair » et « l'ens. » sont écrits au-dessus de la ligne. L'indice de $S_D$ est un $D$ ; son $S$ est un S doublé}
\marginal{Il vaut peut-être mieux dire que $D$ est le support d'un $1$-sous-complexe \ill{} \ill{} du \ill{} $K^{1}$ \ill{} \ill{} \ill{} et ce sous-complexe est un polygône comb.\ d'ordre \ill{} (sans exclure les cas $s_D = 1$ ou $s_D = 2$ \ldots)}
\note{la note marginale est écrite en oblique dans le coin supérieur gauche, sur huit lignes serrées ; « $K^{1}$ » est ajouté au-dessus d'un mot, « d'ordre » au-dessus de la fin d'une ligne}

Soit $s$ un sommet de la carte. Si on exclut le cas où toutes les ps.\ dr.\ sont « concourantes » en $s$ (auquel cas $s$ est l'unique sommet), il existe \uncertain{donc} une $D$ qui ne passe pas par $s$, et l'ens.\ \add{$\Omega_s(X)$} des orientations de $X$ en $s$ est en corr.\ 1-1 can.\ avec l'ens.\ \add{$\Omega(D)$} des orientations du polygône comb.\ $D$ (qui est défini en termes du $1$-complexe $K^{1}$ et de sa décomposition en ps.\ droites \ldots) Mais soient $D'$ \struck{\ill{}} une autre ps.\ droite ne passant pas par $s$, comment expliciter la bijection composée
\[
\Omega(D) \xrightarrow{\ \sim\ } \Omega(D') \qquad \text{(via } \Omega_s(X)\text{)} \ ?
\]
\note{sur la page, $\Omega(D) \xrightarrow{\sim} \Omega(D')$ est suivie d'une flèche oblique $\Omega(D) \xrightarrow{\sim} \Omega_s(X)$ et d'une flèche remontant de $\Omega_s(X)$ vers $\Omega(D')$, marquée $\sim$ ; nous les résumons par « via $\Omega_s(X)$ ». $\Omega_s(X)$ et $\Omega(D)$ sont écrits au-dessus de la ligne}

? \struck{Supposons} Comme $s$ est d'ordre $\geqslant 4$, il passe par $s$ deux ps.\ droites $\Delta$, $\Delta'$, \struck{\ill{}} puis prenons \struck{\ill{}} \ill{} distinctes de $\Delta$ et $\Delta'$, non concourantes \uncertain{sur} $D, D'$ ; \ill{} ; soient $u, u'$ (resp.\ $v, v'$) les int.\ de $\Delta$ (resp.\ $\Delta'$) avec $D$, $D'$, dont l'int.\ est \uncertain{désignée} par $t$. On constate alors (puisque la configuration des quatre ps.\ droites
\note{les deux premières lignes de ce paragraphe sont écrites serrées, l'une au-dessus de l'autre, avec des ajouts interlinéaires ; « deux ps.\ droites » est au-dessus de la ligne. La phrase se poursuit en haut de la page 62}

\page{62}
$D, D', \Delta, \Delta'$ sont en position générale, donc en position standard) que $\overrightarrow{tu}$ et $\overrightarrow{tu'}$ se correspondent. [Plus généralement, si on a les droites $\Delta_1, \ldots, \Delta_k$ passant par $s$ \add{(coupant $D$ en $u_i$ et $D'$ en $u'_i$,} et non par $t$,] et si on prend sur $D \setminus \lbrace t \rbrace$ et $D' \setminus \lbrace t \rbrace$ les relations d'ordre \add{induites par des orientations de $D$, $D'$} qui se correspondent, alors si on ordonne les $\Delta_i$ de façon que les \struck{$D$} $u_i$ aillent en croissant, de même les $u'_i$]. Ceci définit la bijection
\[
\Omega(D) \simeq \Omega(D')
\]
en termes de la combinatoire de $K^{1}$ et de ses ps.\ droites, sauf dans le cas où il n'y a qu'une seule ps.\ droite \add{$\Delta'$} passant par $s$ et non par $t$, (ce qui implique que qu'il existe exactement $2$ ps.\ dr.\ passant par $s$, i.e.\ que l'ordre de $s$ est $4$, et qu'il existe une ps.\ dr.\ $\Delta'$ passant par $s$ et par $t$) Supposons qu'il existe une ps.\ dr.\ $D''$ \add{ne passant pas par $t$,} (distincte de $\Delta'$) \ill{} (ici que la figure ne se réduise à un faisceau de ps.\ dr.\ passant par $t$, et une sécante $\Delta$ de ce faisceau) --- elle est \ill{} $\neq \Delta, \Delta', D, D'$ et ne passe pas par $s$. Soient $x, x'$ les pts d'int.\ de $D''$ avec $D$, $D'$ et $u, u'$ ceux de $\Delta$ avec $D$, $D'$. Si on a $x \neq u$, $x' \neq u'$ alors on peut comparer par ce qui précède les \struck{\ill{}}
\note{les flèches de $\overrightarrow{tu}$ et $\overrightarrow{tu'}$ sont tracées au-dessus des lettres. « (coupant $D$ en $u_i$ et $D'$ en $u'_i$, » et « induites par des orientations de $D$, $D'$ » sont écrits au-dessus de la ligne, dans des bulles ; « ne passant pas par $t$, » de même au-dessus de $D''$. Dans « que qu'il existe » la répétition est sur la page}
\marginal{\emph{NB} $D''$ ne passe pas par $s$ \ill{}, i.e.\ est $\neq \Delta, \Delta'$}

\page{63}
\struck{$\ill{}$} $\Omega(D'')$ \struck{\ill{}} et $\Omega(D)$, et $\Omega(D'')$ et $\Omega(D')$, donc $\Omega(D)$ et $\Omega(D')$, et on gagne. Si par contre toutes les ps.\ droites qui ne passent pas par $t$ passent par $u$, ou si elles passent toutes par $u'$, on est dans le cas de deux faisceaux de ps.\ droites, de sommets $t$ et $u$ (resp.\ $u'$), qui demande encore un traitement à part. Mais supposons qu'il n'existe pas de $D''$ qui ne passe ni par $u$ ni par $u'$ (ni par $t$, sous-entendu) mais qu'il en existe \add{$\Gamma$} qui ne passe pas par $u'$ \add{(mais par $u$)}, et d'autres \add{$\Gamma'$} qui ne passent pas par $u$ (mais par $u'$), alors
\[
\Omega(\Gamma) \simeq \Omega(D') , \qquad \Omega(\Gamma') \simeq \Omega(D)
\]
et on peut déduire de
\[
\Omega(\Gamma) \simeq \Omega(\Gamma')
\]
celui $\Omega(D) \simeq \Omega(D')$) \emph{pourvu que} l'intersection $r$ de $\Gamma$ et $\Gamma'$
\note{devant la première égalité, deux petits traits ; devant la troisième, un signe « $=$ ». « $\Gamma$ » est écrit au-dessus de « qui ne passe », « (mais par $u$) » au-dessus de la fin de ligne. Après « $\Gamma$ et $\Gamma'$ » un large blanc précède la ligne suivante : la fin de la condition n'est pas écrite}

\note{à gauche, un dessin au crayon : deux droites $D$ et $D'$ issues d'un point $t$ (à gauche), une droite verticale $\Delta$ passant par $s$ (en bas), une droite $\Delta'$ issue de $s$ vers la gauche, une droite horizontale épaissie $\Gamma'$ et une droite oblique $\Gamma$ (en haut) ; les points $u$ (sur $D$) et $u'$ (sur $D'$) sont marqués sur la verticale}

par malheur on se trouve encore justement sur $\Delta'$ (cela suppose que $\Gamma$ et $\Gamma'$ forment p.r.\ à $s$ et $t$, $\Delta, \Delta'$ la même figure que \uncertain{forment} $D, D'$) Dans ce cas, les $6$ ps.\ droites $D, D', \Delta, \Delta', \Gamma, \Gamma'$ sont les

\page{64}
$6$ ps.\ dr.\ qui joignent $2$ à $2$ les quatre sommets $u, u', r, t$ de la figure. Mais on constate alors que la bijection \add{\uncertain{$\varphi_s$} :} $\Omega(D) \simeq \Omega(D')$ définie par $s$ n'est autre que celle définie par $r$ (qui est dans la même des deux régions de $X$ découpées par $D, D'$ \ldots)
\note{le signe au-dessus de « $\Omega(D)$ » est \uncertain{lu} $\varphi_s$, suivi de deux points}

\note{un court trait horizontal sépare ce qui précède de ce qui suit}

$K^{1}$, $1$-complexe subdivisé en ps.\ droites (tout sommet appartenant à au moins $2$ ps.\ dr., et $2$ ps.\ droites se rencontrant \add{en un pt et un seul}). Si $D, D'$ sont $2$ ps.\ dr.\ distinctes, et $s$ un sommet qui n'est ni sur $D$, ni sur $D'$, on a une bijection
\[
\varphi_{s;D,D'} : \Omega(D) \simeq \Omega(D')
\]
avec les propriétés suivantes :

a) Sur l'ens.\ des ps.\ droites \add{$D, D', \ldots$} qui ne passent pas par $s$, le syst.\ des $\varphi_{s;D,D'}$ est \uncertain{simplement} transitif.

b) Si $\Delta$ est une ps.\ droite distincte de $D$ et $D'$, coupant $D$ et $D'$ en $x$ et $x'$, \add{(\emph{NB} on n'exclut pas $x = x'$)}, et si $s, s'$ sont deux sommets sur $\Delta$ distincts de $x, x'$, \struck{alors} \struck{\ill{}} on a
\[
\varphi_{s;D,D'} = \varphi_{s';D,D'}
\]
ssi $s, s'$ appartiennent à une même des deux arcs définis par $x$ et $x'$ sur $\Delta$. Si $a, a'$ sont les deux arcs, on peut donc définir $\varphi_{a;D,D'}$ et $\varphi_{a';D,D'}$.
\note{le $\Delta$ de « Si $\Delta$ est » récrit un $D$ souligné. Dans $\varphi_{s;D,D'}$ de la ligne centrée, un petit $D$ est écrit au-dessus de l'indice}

c) Soit $t$ l'\ill{} intersection de $D$ et $D'$, et soient $\Delta_1, \Delta_2$ \add{($i = 1, 2$)} deux ps.\ dr.\ passant par $s$ et non par $t$, \struck{\ill{}} $u_i$ (resp.\ $u'_i$) ($i = 1, 2$) leurs int.\ avec $D$ ($D'$).
\note{la page s'arrête ici ; l'énoncé de c) n'est pas achevé, et la page 65 commence sur autre chose}

\page{65}
syst.\ \add{d'isoms} en question serait \uncertain{unique} par \ill{} \ill{}, \ill{}. Car notons qu'il ne peut y avoir d'unicité quand il existe une ps.\ droite qui est un \emph{bigône} (pas non plus bien sûr s'il y a un monogône, i.e.\ s'il s'agit d'un faisceau de ps.\ dr.) --- mais cela signifie justement qu'il s'agit d'un double faisceau contenant une ps.\ droite qui joint les deux foyers\ldots

On peut maintenant, dans le cas où il n'y a ni monogône ni \emph{bigône}, définir pour tout sommet $s$ $\Omega(s)$ comme l'ens.\ quotient de l'ens.\ des $\Omega(D)$ ($D$ ps.\ dr.\ ne passant pas par $s$) par le syst.\ des isoms $\varphi_{s;D,D'}$. C'est un ens.\ à deux éléments, \struck{\ill{}} \ill{} \ill{} il faut en faire un syst.\ local sur $K^{1}$ i.e.\ pour tout segment $I$ de $K^{1}$ d'extrémités $s, s'$, il faut définir une bijection canonique $\Omega(s) \simeq \Omega(s')$. Par hyp.\ il existe une droite $D$ qui ne passe par $s$ ni $s'$, et par \uncertain{là} on définit $\Omega(s) \simeq \Omega(s')$ comme la composée

\begin{tikzcd}[column sep=small]
\Omega(s) \arrow[rr, "\sim"] \arrow[dr, "\sim"'] & & \Omega(s') \\
& \Omega(D) \arrow[ur] &
\end{tikzcd}

, il résulte de b) que ça ne dépend pas du choix de $D$
\note{« d'isoms » est écrit au-dessus de « syst. », au début de la page ; ce qui précède manque (la page 64 s'arrête sur l'énoncé inachevé de c)). La flèche oblique descendante porte « $\sim$ »}

\page{66}
(c'est équivalent à b)!)

Pour \add{\uncertain{finir}} reconstituer la carte à partir de $K^{1}$, il faudrait encore définir une structure polygonale sur l'ens.\ des arcs issus d'un sommet $s$, et une bijection canonique entre $\Omega(s)$ et l'ens.\ des deux orientations de cette structure polygonale.
\note{un petit astérisque sous le texte ; le reste de la page est blanc}

\page{67}
\note{feuille plus petite, à l'italienne, au crayon}
$\vec{D}, \vec{D}'$ --- ps.\ droites orientées, $\vec{\mathcal{D}}$ leur ens.

$\vec{\mathcal{D}}(s)$ l'ens.\ des ps.\ dr.\ orientées passant par $s$ ($\Longleftrightarrow$ l'ens.\ des arcs d'origine $s$) (cardinal $2$\,\uncertain{$\nu$}$(s)$)

$\rho_{\vec{D},s}$ permutation circulaire de $\vec{\mathcal{D}}(s)$ (si $s \notin |\vec{D}|$).
\note{la lettre rendue $\rho$, ici et jusqu'à la page 70, est un long trait bouclé en haut, lu $\rho$ sans certitude ; $\vec{\mathcal{D}}$ rend son D rond fléché. Dans « cardinal $2\nu(s)$ » la lettre est \uncertain{lue} $\nu$}

a) $\rho_{-\vec{D},s} = \rho_{\vec{D},s}^{-1}$

b) Si $\vec{D}, \vec{D}' \notin \vec{\mathcal{D}}(s)$, alors $\rho_{\vec{D},s}$ et $\rho_{\vec{D}',s}$ sont égales ou \add{\uncertain{opposées}}
\note{le mot sous « égales ou » est écrit en interligne, d'une lecture incertaine}

\struck{c) Si $\rho_{\vec{D},s} = \rho_{\vec{D}',s}$, et $t$ un $2^{\mathrm{e}}$ sommet tel que $t \notin |\vec{D}| \cup |\vec{D}'|$, alors $\rho_{\vec{D},t} = \rho_{\vec{D}',t}$ ssi $s, t$ appartiennent \ill{} \ill{} \ill{}}
\note{l'item c) est barré de sept crochets obliques ; sa dernière ligne, au bas de la feuille, est à demi coupée. Il est repris à la page suivante}

\page{68}
c) Soient $\vec{D}, \vec{D}', \vec{\Delta}$ ps.\ droites \add{(\uncertain{orientées} \uncertain{supp.}\ distincts)}, $x, x'$ les int.\ de $\Delta$ avec $D, D'$ (on ne suppose pas $x \neq x'$) et $s, s'$ des sommets sur $\Delta$ distincts de $x, x'$. Supposons $\rho_{\vec{D},s} = \rho_{\vec{D}',s}$. Pour qu'on ait $\rho_{\vec{D},s'} = \rho_{\vec{D}',s'}$, il faut et suffit que $s, s'$ appartiennent à un même arc de $\Delta$ déterminé par $x, x'$.

(d) Soit $\vec{\Delta}_1 \in \vec{\mathcal{D}}(s)$, $\vec{D} \in \vec{\mathcal{D}} \setminus \vec{\mathcal{D}}(s)$, $\vec{\Delta}_2 = \rho_{\vec{D},s}(\vec{\Delta}_1)$, soient $\lbrace u_1 \rbrace = |\vec{\Delta}_1| \cap |\vec{D}|$, $\lbrace u_2 \rbrace = |\vec{\Delta}_2| \cap |\vec{D}|$, alors $u_2$ est le sommet sur $\vec{D}$ qui suit $u_1$, \emph{parmi} toutes les intersections de $\vec{D}$ avec les ps.\ dr.\ passant par $s$.
\note{« (d) » : la lettre est entourée. $\vec{\Delta}_2$ porte en exposant un petit signe \ill{} ; l'argument de $\rho_{\vec{D},s}$ est écrit $\vec{\Delta}$ sans indice. « parmi » est souligné deux fois}

e)
\note{la page s'arrête sur ce « e) » ; l'item est repris en tête de la page 69}

\page{69}
\note{à gauche, un dessin : deux ps.\ droites orientées $\vec{D}$ et $\vec{D}'$ issues d'un point $t$, coupées par deux droites $\Delta_1$ et $\Delta_2$ qui se rencontrent en $s$ (en bas) ; sur $\vec{D}$ les points $u_1$, $u_2$, sur $\vec{D}'$ les points $u'_1$, $u'_2$ ; une flèche courbe va de $\Delta_1$ à $\Delta_2$ près de $s$, et $\Delta_2$ porte le même petit exposant qu'à la page 68}
(e) Soient $\vec{D}, \vec{D}' \in \vec{\mathcal{D}} \setminus \vec{\mathcal{D}}(s)$, \struck{$\Delta$} $\lbrace t \rbrace = |\vec{D}| \cap |\vec{D}'|$ (on suppose $|\mathcal{D}| \neq |\mathcal{D}'|$ \add{dist.}) $\Delta_1, \Delta_2$ deux ps.\ droites passant par $s$ et non par $t$, $u_1, u_2$ ($u'_1, u'_2$) leurs int.\ avec $\vec{D}$, $\vec{D}'$. 1°) Supposons que $u_1$ et $u_2$ soient adj.\ sur $\vec{D}$, alors $u'_1$ et $u'_2$ sont adj.\ sur $\vec{D}'$. Supposons $\vec{D}$ orientée dans le sens $t u_1 u_2$, $\vec{D}'$ dans le sens $t u'_1 u'_2$, et que $u_1, u_2$ soient adj.\ dans l'ens.\ des \struck{ps.\ droites} \add{int.\ de $\vec{D}$} avec les ps.\ dr.\ qui passent par \uncertain{$\vec{D}$}. Alors \struck{$\rho(s u_1 \ldots$}
\[
\rho_{\vec{D}}(s u_1 u'_1) = (s u_2 u'_2) .
\]
\note{« (e) » est entouré. Dans « $|\mathcal{D}| \neq |\mathcal{D}'|$ » les D sont ronds, sans flèche. « qui passent par $\vec{D}$ » est écrit tel quel ; on attendrait « par $s$ », comme en (d). L'indice de $\rho_{\vec{D}}$ ne porte pas de $s$ ; au-dessus, « $\vec{D}$ » récrit une première écriture biffée}

(f) Si $s$ est d'ordre $2n$, alors
\[
(\rho_{\vec{D},s})^{n}(\vec{\Delta}) = -\vec{\Delta} .
\]
\note{« (f) » est entouré}

\page{70}
\emph{Th} Si un syst.\ de $\rho_{\vec{D},s}$ existe satisfaisant a) b) c) d) e) f) \emph{et si} il n'existe pas de ps.\ droites qui sont des monogônes ou bigônes, alors ce système des $\rho_{\vec{D},s}$ est unique. De plus, le syst.\ des $\Omega(s) = \lbrace \rho_s, \rho_s^{-1} \rbrace$ --- en \ill{} (\ill{} \ill{}) de \emph{façon unique} un syst.\ local satisfaisant \struck{\ill{}} \ldots
\note{le second élément de $\lbrace \rho_s, \rho_s^{-1} \rbrace$ est écrit $\rho''_s$ ou $\rho_s^{-1}$, \uncertain{lu} $\rho_s^{-1}$. Après « satisfaisant » un mot est noirci d'un gribouillis en zigzag, puis quatre points}
\marginal{il suffit de d) e), f) et (ss doute) a) en est une conséquence. Pour pouvoir donner un sens au syst.\ local, il faut des \ill{} plus \emph{fortes} \uncertain{que} \uncertain{les} hyp.\ \uncertain{suppl.}\ b) et c) \struck{pour pouvoir \ill{}}}
\note{la note marginale est écrite en oblique dans la marge gauche, sur une dizaine de lignes courtes}

\page{72}
\emph{1 droite} : standard pr.

\emph{2 droites} : standard pr.\ autom.\ proj/comp.\ neutre $\simeq$ autom.\ comb.\ $\simeq$ aut.\ \ill{}

\emph{3 droites}

a) \struck{position} générale : standard pr.\ endom proj/comp.\ neutre $\simeq$ autom.\ comb.\ $\simeq$ aut.\ \ill{} \add{\uncertain{orientés}} $\simeq$ autom.\ de l'ens.\ des \uncertain{quatre} faces

b) concourantes : standard proj ? end.\ proj/comp.\ neutre $\simeq$ \add{ext.\ de} \struck{\ill{}} $\mathfrak{S}_3$ \ldots\ $\mathbf{Z}/2$ ;
\struck{\ill{}} $\mathfrak{S}_3$ $\simeq$ autom.\ comb.\ $\simeq$ ext.\ de $\mathbf{Z}/2$ par $\mathbf{Z}/6$ $\simeq$ groupe de l'hexagone (?)
\note{les deux items sont réunis par une accolade ; une longue courbe sépare, à droite, la colonne de a) de celle de b). « ext.\ de » est écrit au-dessus d'un mot biffé suivi d'un $\mathfrak{S}_3$ entouré, puis de $\mathbf{Z}/2$ ; le signe devant le second $\mathfrak{S}_3$ est noirci, et un petit signe illisible est écrit au-dessus}

\emph{4 droites}

a) pos.\ gén.\ : standard proj.\ [endom proj $\simeq$ Id \add{mais} endom comb.\ $\simeq \mathfrak{S}_4$

b) \struck{trois} droites conc.\ : standard proj., endom proj/comp.\ $\simeq$ ext.\ de $\mathfrak{S}_3$ par $\mathbf{Z}/2$ $\simeq$ \struck{ext.} autom.\ comb.\ $\simeq$ gr.\ hexagone

c) quatre droites conc.\ [pas standard projectif \struck{\ill{} autogonal} \add{mais std comb}] endom proj : \ill{} mais autom comb.\ groupe de l'octogone (?)
\note{les trois items sont réunis par une accolade. En c), « pas standard projectif » et « mais std comb » sont encadrés, et la ligne biffée entre les deux est d'une lecture incertaine}

\emph{5 droites} a) position générale : \struck{pas} standard proj., mais standard comb., \struck{groupe des autom comb $=$} groupe du pentagone

$n \geqslant 6$ \emph{droites} pos.\ générale : \emph{pas} standard \emph{comb.}
\note{« 1 droite », « 2 droites », etc.\ sont soulignés, et forment une colonne à gauche ; « pas » et « comb. » de la dernière ligne sont soulignés}

\end{document}
